\documentclass[12pt,a4paper]{article}

\usepackage[utf8]{inputenc}
\usepackage[T1]{fontenc}
\usepackage{amsmath, amssymb, amsthm}
\usepackage{graphicx}

\title{\textbf{Die $\vec{\Omega}$-Theorie der Gravitation} \\ \large Eine rotationsbasierte, nicht-lokale Feldtheorie, emergiert aus der WDBT+}
\author{Michael Czybor}
\date{März 2026}

\newtheorem{axiom}{Axiom}
\newtheorem{definition}{Definition}
\newtheorem{satz}{Satz}

\begin{document}

\maketitle

\begin{abstract}
Die $\vec{\Omega}$-Theorie ist eine nicht-lokale, instantane Feldtheorie der Gravitation. Sie basiert auf der fundamentalen Idee, dass Gravitation nicht durch Krümmung der Raumzeit, sondern durch ein fundamentales Rotationsfeld $\vec{\Omega}$ vermittelt wird. Die Theorie emergiert aus der feldlosen WDBT+ (Weber-de Broglie-Bohm-Theorie mit fraktalem Raum) und der DSTT (Dynamischen Schwere-Trägheits-Theorie). Sie stellt die Rotationsanalogie zur Allgemeinen Relativitätstheorie dar, unterscheidet sich jedoch fundamental durch ihre Nicht-Lokalität, Instantaneität und die intrinsische Irreversibilität ($\dot{\beta} > 0$), die zur Nicht-Erhaltung des Drehimpulses führt.
\end{abstract}

\tableofcontents
\newpage

\section{Einleitung: Vom Teilchen zum Feld}

Die WDBT+ beschreibt die Gravitation als direkte, instantane Teilchen-Teilchen-Wechselwirkung nach dem Weber-Gesetz, gesteuert durch das Bohmsche Quantenpotential. Sie ist fundamental \textit{feldlos}. Die DSTT übersetzt diese Dynamik in die Sprache der Energieaufteilung zwischen radialer und tangentialer Komponente, eingefangen in der Zustandsvariablen $\beta(t)$ mit der fundamentalen Eigenschaft $\dot{\beta} > 0$.

Die $\vec{\Omega}$-Theorie ist der nächste Schritt der Emergenz. Sie ist der \textit{Kontinuumslimes} der DSTT und damit die effektive Feldtheorie für makroskopische Systeme. An die Stelle der direkten Teilchen-Wechselwirkung treten Felder, die die kollektive Wirkung vieler Massenpunkte beschreiben. Das zentrale neue Element ist das Vektorfeld $\vec{\Omega}$, das die lokale Rotation des Raumes selbst repräsentiert. Die $\vec{\Omega}$-Theorie verhält sich zur ART wie die Maxwell-Theorie zur Newton-Mechanik, mit dem entscheidenden Unterschied, dass sie die Nicht-Lokalität und Irreversibilität ihrer fundamentalen Basis bewahrt.

\section{Axiome der $\vec{\Omega}$-Theorie}

\begin{axiom}[Nicht-Lokalität und Instantaneität]
Die Gravitationswirkung ist nicht-lokal und instantan. Es gibt keine Ausbreitung mit endlicher Geschwindigkeit; die Wechselwirkung ist simultan. Dieser Axiom ist direkte Konsequenz der feldlosen Weber-Wechselwirkung der WDBT+.
\end{axiom}

\begin{axiom}[Das Rotationsfeld]
Die Gravitation wird durch ein fundamentales, axiales Vektorfeld $\vec{\Omega}(\vec{r}, t)$ beschrieben. Dieses Feld kodiert die lokale Rotationsdichte des Raumes und ist die Ursache für alle geschwindigkeitsabhängigen (tangentialen) Trägheitseffekte.
\end{axiom}

\begin{axiom}[Aufspaltung von Ursache und Wirkung]
Die Schwere (Ursache) und die Trägheit (Wirkung) sind fundamental verschieden. Die radiale Anziehung wird durch ein skalares Potential $\Phi(\vec{r}, t)$ beschrieben, die tangentialen Trägheitskräfte durch das Rotationsfeld $\vec{\Omega}(\vec{r}, t)$.
\end{axiom}

\begin{axiom}[DSTT-Kompatibilität]
Die Theorie ist so konstruiert, dass aus ihr die Existenz einer Zustandsvariablen $\beta(\vec{r}, t)$ folgt, die das lokale Verhältnis von tangentialer zu gesamter Trägheit angibt, und für die $\dot{\beta} > 0$ gilt. Dies verleiht der Gravitation einen intrinsischen Zeitpfeil.
\end{axiom}

\section{Die Feldgleichungen}

Die Felder werden nicht-lokal und instantan aus den Quellen erzeugt.

\subsection{Definition der Felder}

\begin{definition}[Skalares Potential]
Das Potential der radialen Anziehung (gravito-elektrisches Potential) ist durch die Massendichte $\rho$ bestimmt:
\begin{equation}
\boxed{ \Phi(\vec{r}, t) = -G \int_{\vec{r}' \neq \vec{r}} d^3r' \, \frac{\rho(\vec{r}', t)}{|\vec{r} - \vec{r}'|} }
\end{equation}
Der Ausschluss der Selbstwechselwirkung ($\vec{r}' \neq \vec{r}$) ist essentiell für die Konsistenz der Theorie und vermeidet Divergenzen.
\end{definition}

\begin{definition}[Vektorpotential]
Das Vektorpotential, aus dem die Rotation abgeleitet wird, wird durch die Massenstromdichte $\vec{j}(\vec{r}, t) = \rho(\vec{r}, t) \vec{v}(\vec{r}, t)$ erzeugt:
\begin{equation}
\boxed{ \vec{A}(\vec{r}, t) = \frac{G}{c^2} \int_{\vec{r}' \neq \vec{r}} d^3r' \, \frac{\vec{j}(\vec{r}', t)}{|\vec{r} - \vec{r}'|} }
\end{equation}
\end{definition}

\begin{definition}[Rotationsfeld]
Das fundamentale Rotationsfeld $\vec{\Omega}$ ist die Wirbelstärke des Vektorpotentials:
\begin{equation}
\boxed{ \vec{\Omega}(\vec{r}, t) = \nabla \times \vec{A}(\vec{r}, t) }
\end{equation}
\end{definition}

\begin{definition}[Induziertes Feld]
Bei zeitabhängigen Strömen entsteht ein induziertes radiales Feld:
\begin{equation}
\boxed{ \vec{E}_{\text{ind}}(\vec{r}, t) = -\frac{\partial \vec{A}(\vec{r}, t)}{\partial t} }
\end{equation}
\end{definition}

\subsection{Die Maxwell-Form der Feldgleichungen}

Die so definierten Felder gehorchen im Vakuum (außerhalb der Quellen) Gleichungen, die formal den Maxwell-Gleichungen der Elektrodynamik entsprechen:

\begin{align}
\nabla \cdot \vec{\Omega} &= 0 \label{eq:max1} \\
\nabla \times \vec{\Omega} &= \frac{4\pi G}{c^2} \, \vec{j} + \frac{1}{c^2} \frac{\partial}{\partial t} (-\nabla \Phi) \label{eq:max2} \\
\nabla \cdot (-\nabla \Phi) &= -4\pi G \rho \label{eq:max3} \\
\nabla \times (-\nabla \Phi) &= -\frac{\partial \vec{\Omega}}{\partial t} \label{eq:max4}
\end{align}

Es ist jedoch entscheidend zu betonen, dass diese Differentialgleichungen hier nicht fundamental sind, sondern eine alternative Schreibweise der nicht-lokalen Integralgleichungen (1) und (2) darstellen. Die wahre Natur der Theorie ist die der instantanen, integralen Wechselwirkung.

\section{Die Bewegungsgleichung}

Eine Probemasse $m$ am Ort $\vec{r}$ mit der Geschwindigkeit $\vec{v}$ erfährt die Kraft:
\begin{equation}
\boxed{ \vec{F} = m \left( -\nabla \Phi + 2 \, \vec{v} \times \vec{\Omega} - \frac{\partial \vec{A}}{\partial t} \right) }
\end{equation}
Dies ist die Lorentz-Kraft der Gravitation. Der Term $-\nabla \Phi$ repräsentiert die radiale Newtonsche Anziehung. Der Term $2\vec{v} \times \vec{\Omega}$ ist die geschwindigkeitsabhängige Coriolis-artige Kraft, die für die tangentialen Trägheitseffekte verantwortlich ist. Der Term $-\partial \vec{A}/\partial t$ beschreibt induzierte Effekte durch instationäre Massenströme.

\section{Die DSTT-Variable $\beta$}

In der $\vec{\Omega}$-Theorie lässt sich die DSTT-Zustandsvariable $\beta$ direkt als lokales Verhältnis der Kraftkomponenten definieren:
\begin{equation}
\boxed{ \beta(\vec{r}, t) = \frac{ |\vec{v} \times \vec{\Omega}| }{ |\nabla \Phi| + |\vec{v} \times \vec{\Omega}| } }
\end{equation}

Aus den Feldgleichungen und der Bewegungsgleichung folgt die für die DSTT fundamentale Eigenschaft:
\begin{equation}
\boxed{ \dot{\beta} > 0 }
\end{equation}
Dies ist keine zusätzliche Annahme, sondern eine Konsequenz der nicht-lokalen, instantanen Rückkopplungsdynamik der Theorie. Die zeitliche Änderung von $\vec{\Omega}$ und $\vec{v}$ führt zu einer monotonen Zunahme des tangentialen Anteils der Bewegung, was die Bahnen zu offenen Spiralen macht.

\section{Zentralfeld-Lösung und Spiralbahnen}

Für das Zweikörperproblem (Zentralmasse $M$, Testmasse $m$) liefert die $\vec{\Omega}$-Theorie eine Bewegungsgleichung, die auf eine logarithmische Spirale führt. Mit dem Ansatz $\Omega_\theta = k/r^2$ (der sich aus den Feldgleichungen ergibt) folgt aus den Gleichungen in Polarkoordinaten:
\begin{align}
\ddot{r} - r\dot{\phi}^2 &= -\frac{GM}{r^2} + 2 r\dot{\phi} \, \Omega_\theta \\
r\ddot{\phi} + 2\dot{r}\dot{\phi} &= -2\dot{r} \Omega_\theta
\end{align}
Die Lösung dieser Gleichungen ist:
\begin{equation}
\boxed{ r(\phi) = K e^{\phi / B} }
\end{equation}
wobei $K$ und $B$ Konstanten sind, die von den Anfangsbedingungen und der Stärke des $\vec{\Omega}$-Feldes abhängen. Es gibt keine geschlossenen Bahnen (Ellipsen) – ein direktes Observatorium der DSTT und der Nicht-Erhaltung des Drehimpulses.

\section{Kosmologische Konsequenzen}

Die $\vec{\Omega}$-Theorie führt zu einem stationären, nicht-expandierenden Universum.
\begin{itemize}
    \item \textbf{Rotverschiebung:} Sie ist keine Doppler-Folge der Expansion, sondern eine kumulative gravitative Wirkung auf Photonen während ihrer Reise durch das $\vec{\Omega}$-Feld.
    \item \textbf{Galaxienrotation:} Flache Rotationskurven werden ohne Dunkle Materie durch das $\vec{\Omega}$-Feld erklärt.
    \item \textbf{CMB-Dipol:} Der beobachtete Dipol im Mikrowellenhintergrund ist eine Superposition aus der lokalen Bewegung des Sonnensystems und einem intrinsischen Anteil, der von einem kosmologischen $\vec{\Omega}_{\text{kosmisch}}$-Feld herrührt.
    \item \textbf{Zeitpfeil:} Die Bedingung $\dot{\beta} > 0$ liefert einen intrinsischen kosmologischen Zeitpfeil, der nicht aus der Thermodynamik abgeleitet werden muss.
\end{itemize}

\section{Zusammenfassung: Die $\vec{\Omega}$-Theorie im Kontext}

\begin{table}[ht]
\centering
\begin{tabular}{|l|c|c|}
\hline
\textbf{Eigenschaft} & \textbf{ART} & \textbf{$\vec{\Omega}$-Theorie} \\
\hline
Fundament & Raumzeit-Krümmung & Rotationsfeld $\vec{\Omega}$ \\
Lokalität & Ja (Differentialgleichungen) & \textbf{Nein} (Integralgleichungen) \\
Ausbreitung & Mit Lichtgeschwindigkeit & \textbf{Instantan} \\
Drehimpuls & Erhalten & \textbf{Nicht erhalten} ($\dot{L} \neq 0$) \\
Intrinsischer Zeitpfeil & Nein & \textbf{Ja} ($\dot{\beta} > 0$) \\
DSTT-Kompatibel & Nein & \textbf{Ja} \\
Fundamentalität & Fundamental & \textbf{Emergent} (aus WDBT+) \\
\hline
\end{tabular}
\caption{Vergleich zwischen der Allgemeinen Relativitätstheorie und der $\vec{\Omega}$-Theorie.}
\end{table}

Die $\vec{\Omega}$-Theorie ist die natürliche Feldtheorie der Gravitation, die aus der feldlosen, teilchenbasierten WDBT+ emergiert. Sie stellt eine vollwertige, rotationsbasierte Alternative zur ART dar und erweitert diese um die essentiellen Konzepte der Nicht-Lokalität, Instantaneität und Irreversibilität, die zur Erklärung kosmologischer Phänomene ohne Dunkle Materie oder Dunkle Energie notwendig sind.

\end{document}